\section{Постановка задачи}

Необходимо выполнить пять практических работ по функциональному программированию на языке Haskell. Работы должны быть размещены в репозитории по отдельным каталогам, изменения должны фиксироваться в Git, а для проектов со сборкой должны использоваться Cabal или Stack.

В практической работе №1 требуется реализовать набор функций для работы со списками, типом \mintinline{haskell}{Maybe}, пользовательскими типами данных и функциями приготовления тортов из лекционных материалов. На защите требуется решить задачу построения функции типа \mintinline{haskell}{Either (a -> a) (c -> c) -> [a] -> [c] -> [(a, c)]}.

В практической работе №2 требуется реализовать пользовательские типы данных, сделать их представителями классов типов, написать парсеры и разобрать выражение с applicative-оператором \mintinline{haskell}{(<*>)} для типа \mintinline{haskell}{(Either a Int, String)}.

В практической работе №3 требуется реализовать прикладные проекты: обработку изображения с последовательными глитч-операциями и поиск пути в лабиринте. Для лабиринта необходимо использовать монаду \mintinline{haskell}{RWS} и объяснить, почему тип функции поиска строится через \mintinline{haskell}{RWS}.

В практической работе №4 требуется создать Stack-проект текстовой игры по лабиринту, прочитать лабиринт из файла, реализовать собственный трансформер состояния \mintinline{haskell}{MyStateT}, собрать стек трансформеров и реализовать запуск игры. На защите необходимо объяснить запуск трансформера, порядок терминирования и места вызова \mintinline{haskell}{IO}.

В практической работе №5 требуется создать Stack-проект \mintinline{text}{qc}, реализовать функции \mintinline{haskell}{addMod} и \mintinline{haskell}{reverseWords}, написать генераторы входных данных и проверить свойства функций с помощью QuickCheck.

Сдачи практических работ:

\begin{itemize}
    \item 20.03.2026 --- сдана практическая работа №1;
    \item 24.04.2026 --- сдана практическая работа №2;
    \item 22.05.2026 --- сдана практическая работа №3;
    \item 22.05.2026 --- по практической работе №4 остался один вопрос;
    \item практическая работа №5 подготовлена в репозитории.
\end{itemize}

\newpage
